\section{Controlled identification protocols}
\label{subsec:synthetic-protocol-details}

The planned evaluation tests nonlinear action-block identification,
controlled two-block recombination, and rendered-video action transfer.
The numerical generators have known latent variables and checked
observation properties. Finite-data training tests the predicted recovery
behavior; it does not establish the population theorem experimentally.

\paragraph{Nonlinear action-anchor identification.}
Sample a complete initial-state, control, and innovation realization,
then independently resample appearance to construct
Eq.~\ref{eq:controlled-anchor-views}. The primary process is
Proposition~\ref{prop:nonlinear-temporal-example}, using the shear map
in Eq.~\ref{eq:nonlinear-shear-observation} and additional invertible
nonlinear mixing compositions. Three or more distinct times meet the
main conditions; one- and two-time observations provide controlled
failures of temporal observability. The three trajectory parameters
specify the anchor dimension without supervising its coordinates.

Appendix~\ref{subsec:numerical-verification} trains agreement plus kernel
distribution prediction, using Eq.~\ref{eq:empirical-kernel-prediction},
and compares it with ordinary-mean prediction, agreement alone, and
prediction alone. For continuous data, the kernel dictionary and neural
encoder approximate the population class; held-out trajectories test this
approximation, and increasing their capacity remains a further test. Exact finite-support
calculations provide an additional check of the objective's minimizers.
The multiplication model in Eq.~\ref{eq:nonlinear-mean-counterexample}
explicitly tests whether the learner retains action information that
ordinary means erase. Ground-truth factors and mixing inverses are
reserved for mathematical checks and evaluation, not training targets.

Held-out trajectory recovery is reported jointly with within-action
variation and the change in code scale during training.
Temporal pooling and shared appearance across views test which information
or interventions are lost; reversal controls and restricted appearance
variation remain planned.
These numerical identification studies are separate from the patch-view
Action-VAE and rendered-video evaluations.

\paragraph{Additive special case.}
The two constructions in Appendix~\ref{subsec:temporal-observability}
use nonlinear dynamics with an appearance offset, and constant-acceleration
dynamics with invertible linear mixing. They test
Condition~\ref{cond:additive-observations}, the covariance projection in
Proposition~\ref{prop:structural-anchor-recovery}, and ordinary-feature
learning under Corollary~\ref{cor:additive-anchor-learning}.
The specified encoder and readout classes contain the constructed maps.
These compact continuous constructions inherit the main model's sampling
and regularity, and their temporal separation verifies its observability.
They therefore test a subclass of the nonlinear theory and the further
simplification to ordinary-feature prediction. Singular content covariance
and fixed offsets remain compatible with full content support and are not
labeled as violations. Appendix~\ref{subsec:numerical-verification} reports
the covariance projection and ordinary-feature learning for both constructions.

\paragraph{Complementary-intervention model.}
The experiment will independently sample continuous content factors and
action-state trajectories with full support over their Cartesian product,
generate both intervention types in Eq.~\ref{eq:dgp-controlled-pairs}, and
use a known injective observation map.
One construction is a composition of square full-rank affine maps and
LeakyReLU activations with strictly positive negative slopes, applied to
the concatenated content and trajectory variables.
This gives an explicitly invertible numerical observation model.
The content and action coordinate blocks of $g_{\mathrm{ctrl}}^{-1}$,
together with decoder $g_{\mathrm{ctrl}}$, form a zero-loss solution to
Eq.~\ref{eq:identification-objective}. The learner will use matching
block dimensions and function classes containing these maps.
It will expose separate content and temporal action codes and reconstruct
through those codes without a raw-input bypass.
The numerical construction directly tests block identification; the
rendered-video arm evaluates transfer using the simulator's physical
trajectories and counterfactual outputs.
Ground-truth values are reserved for evaluation in the main
intervention-supervised arm; the simulator supplies pair membership.
A separately labeled truth-supervised arm will measure the benefit of
coordinate supervision and connect to the existing synthetic Action-VAE.
An unrestricted reconstruction baseline and removal of each invariance
term will test the identification conditions. Removal of temporal
regularization will measure its contribution to finite-sample recovery
and transfer.

\paragraph{Recovery and leakage.}
For each continuous ground-truth factor $y_j$, a readout $r_j$ will be
fitted on training codes, selected on validation groups, and frozen
before testing. Write $\mathbf e_i$ for the code of sample $i$, $y_{ij}$ for its
$j$-th factor, and $\bar y_{\mathrm{test},j}$ for that factor's mean over
the test index set $\mathcal I_{\mathrm{test}}$. Its test score is
\begin{equation}
  R_j^2
  =1-\frac{\sum_{i\in\mathcal I_{\mathrm{test}}}
                (y_{ij}-r_j(\mathbf e_i))^2}
              {\sum_{i\in\mathcal I_{\mathrm{test}}}
                (y_{ij}-\bar y_{\mathrm{test},j})^2}.
  \label{eq:recovery-r2}
\end{equation}
Negative scores will be retained. Factors with zero test variance will
be reported as undefined, with their coverage and absolute error,
rather than assigned a successful score. We will report per-factor
scores and macro averages over prespecified nonredundant groups.

Because blockwise identification permits nonlinear transformations within each
block, both a linear ridge probe and a nonlinear kernel ridge probe
will be reported. Nonlinear readouts for block recovery have precedent
in invariance-based identification studies~\citep{kugelgen2021content}.
A single readout will be shared across test scenes; test-specific
alignment and fitting on test labels are excluded. For a trajectory
block, the readout may consume the full ordered code sequence.
Phasewise linear decoding is a stricter accessibility diagnostic.

With $\widehat{\mathbf z}^{\mathrm c}=E_{\mathrm c}(V)$ and $\widehat{\mathbf z}^{\mathrm a}=E_{\mathrm a}(V)$ denoting the two controlled-model
codes, we will evaluate all four block-to-factor mappings:
\begin{equation}
  \renewcommand{\arraystretch}{1.5}
  \begin{array}{c|cc}
    \text{readout input} & \text{content target } \mathbf z^{\mathrm c}
                         & \text{action target } \mathbf z^{\mathrm a}_{0:T}\\ \hline
    \widehat{\mathbf z}^{\mathrm c} & R^2_{\mathbf z^{\mathrm c}\leftarrow \widehat{\mathbf z}^{\mathrm c}} & R^2_{\mathbf z^{\mathrm a}\leftarrow \widehat{\mathbf z}^{\mathrm c}}\\
    \widehat{\mathbf z}^{\mathrm a} & R^2_{\mathbf z^{\mathrm c}\leftarrow \widehat{\mathbf z}^{\mathrm a}} & R^2_{\mathbf z^{\mathrm a}\leftarrow \widehat{\mathbf z}^{\mathrm a}}
  \end{array}.
  \label{eq:recovery-leakage-matrix}
\end{equation}
With the independent content/action sampling of the controlled experiment,
desired-block recovery should be high and cross-block prediction should
approach its appropriate constant-predictor baseline. High action
$R^2$ alone does not establish disentanglement: a code containing both
content and action can also attain it. Low probe performance is a
diagnostic, not a proof that no information is present.
For datasets with correlated content and action, leakage must instead
be assessed under controlled interventions or relative to a conditional
baseline.

\paragraph{Semantic and temporal metrics.}
The existing 48-dimensional truth-supervised model will additionally be
evaluated by direct coordinate $R^2$, without fitting a corrective
readout, for its continuous channels. Derived quantities such as
velocity and acceleration will be reported by group rather than
counted as independent causes. Contact states and discrete events will
use macro F1, balanced accuracy, and event-timing error where applicable.
Phase-only and instruction-only baselines will test whether fixed
timing or action labels explain recovery without video evidence.
Recovery will be measured before packet centering and RMS normalization,
which can remove offsets and amplitude; the normalized packet will
receive a separate action-transfer evaluation.

MCC after a fixed, validation-selected one-to-one coordinate matching
can be reported as an additional diagnostic, using rank correlation
when appropriate~\citep{yao2022leap}. It is not the primary test of a
block-level theorem, which permits within-block mixing.
A componentwise identification claim would require a separate theorem
and matching synthetic assumptions. A graph-recovery metric is relevant
only if a graph is explicitly estimated and compared with generative
ground truth.

\paragraph{The implemented representation.}
For a fixed physical rollout, we will rerender independent appearance
variants while holding dynamics, geometry, and camera fixed.
This supplies controlled appearance interventions; a simulator renderer
must also retain the evaluated action information for exact identification.
Frozen features need not have an additive decomposition.
Additional camera variation will be evaluated separately.
We will train the specified Action-VAE with its patch views and losses,
then measure action-descriptor recovery and content dependence for
the raw mean code, centered code, normalized code, projected V-JEPA2
branch, and their complete concatenation.
Descriptors include relative displacement, direction, and interaction
order and are fixed before fitting probes; they are not redefined to
match information discarded by normalization.
The 48-channel supervised model is a separately labeled comparison.

\paragraph{Information retention and content dependence.}
A single training-set readout will be evaluated on held-out trajectories
and appearance combinations. Both linear and nonlinear probes will be
reported; a failed probe does not establish absence of information.
Controlled appearance pairs directly measure
Eq.~\ref{eq:conditioning-content}, with branch scales fixed from training
data. Within-action stability is reported jointly with recovery,
preventing a constant representation from appearing successful.
Nonlinear appearance--motion mixing is included within the theorem-aligned
study when temporal observability holds. Shared-content pairs,
non-observable windows, and observation noise test departures from its
premises. Singular content covariance and fixed offsets separately test
the additive construction within its sufficient conditions.

\paragraph{Editing with the complete condition.}
Matched-budget ablations will compare the action branch, projected V-JEPA2
branch, and their concatenation, with separate models trained for each
input configuration. Retaining the code's removed mean and RMS is an
additional proposed ablation, not part of the current method.
Tests fix source, instruction, and seed while changing demonstration
appearance or its direction and temporal order.
Final-output behavior and source fidelity assess execution; intermediate
reader agreement is reported separately.
Recovery will be conditioned on source and instruction where they
already determine part of the requested behavior.

\paragraph{Transfer, splits, and robustness.}
For synthetic videos, source-content/reference-action swaps will be
compared with the simulator's counterfactual output using action-state
error, content preservation, and event execution. Reference appearance
will vary while the realized action remains fixed.
Zero, reversed, shuffled, and wrong-action references will test whether
the editor uses the ordered action signal. The MuJoCo
simulator supplies the physical rollouts and counterfactual outputs for
this evaluation.

Train, validation, and test partitions will be assigned to independent
base scenes and trajectory realizations before frames, views, or
counterfactual pairs are expanded. No trajectory or render variant
will cross these partitions. A separate compositional split may reuse
individual factor values while withholding content--action
combinations; it will be labeled distinctly.
Scalers, probe settings, and coordinate assignments will use training
and validation data only. Results will include multiple generation and
training seeds, with uncertainty aggregated over scenes or complete
trajectories rather than treating frames as independent samples.

Finally, controlled violations will remove one intervention type,
restrict action coverage, introduce non-injective observations, or add
camera and content-dependent dynamics. Independent observation noise also tests the additive covariance
construction; an estimated low-variance subspace is reported separately
from its noiseless population formula. Nonlinear content-dependent
observation mixing remains within the main theorem when the specified
observability and regularity conditions hold, although it can violate
the additive subclass. Each stress test will state which premise it changes. Real-world
experiments, where ground-truth latents are unavailable and latent
$R^2$ cannot be computed, instead assess edited-video behavior and
source fidelity under uncontrolled departures from these assumptions.

\subsection{Numerical verification of action identification}
\label{subsec:numerical-verification}

We simulate the constructions of
Appendices~\ref{subsec:nonlinear-temporal-example}
and~\ref{subsec:temporal-observability} with known latent variables.
A video is an ordered array of low-dimensional frames. Training uses only
action-preserving pairs; latent values are used only for evaluation.
These finite-data results test the predicted recovery behavior; they do
not establish the population theorem experimentally.

\paragraph{Data.}
The action is $\boldsymbol\theta=(p,v,u)$, the position, velocity, and
constant acceleration at the midpoint of a fixed window, drawn uniformly
from $[-1,1]^3$. Frame $t$ observes $p_t=p+vs_t+\tfrac12us_t^2$ at $T$
evenly spaced times $s_t\in[-1,1]$, with $s_1=0$ for $T=1$.
This affine reparametrization of $(p_0,v_0,u)$ in
Proposition~\ref{prop:nonlinear-temporal-example} leaves its observability
argument unchanged, and the symmetric window maps time reversal to
$(p,-v,u)$ within the support. Frames are $\mathbf x_t=\rho(c,p_t/2.5)$
for two mixing maps: the shear map in
Eq.~\ref{eq:nonlinear-shear-observation} with $c\sim\operatorname{Unif}[-1,1]$,
and the multiplicative map $\rho(c,p)=(cp,c)$ with $c$ uniform on
$[-2,-1]\cup[1,2]$. The latter extends
Eq.~\ref{eq:nonlinear-mean-counterexample} to trajectories: it is
invertible on its support, yet $\E[V\mid\mathbf z^{\mathrm a}]=\mathbf 0$.
Each pair shares $\boldsymbol\theta$ and draws its two contents
independently (Eq.~\ref{eq:controlled-anchor-views}). We use 20{,}000
training pairs, and 4{,}000 validation and 4{,}000 test pairs with new
trajectories.

\paragraph{Learning and evaluation.}
The encoder is a multilayer perceptron with three hidden layers of width
256 applied to the vectorized ordered video, with $d_{\mathrm a}=3$.
The predictor has the form of Eq.~\ref{eq:empirical-kernel-prediction},
with 1{,}024 training videos as its dictionary and a multilayer perceptron
for $\boldsymbol\alpha_\theta$. The kernel averages Gaussian kernels with
bandwidths of $0.25$, $0.5$, and $1$ times the median distance between
dictionary videos; the average is bounded, continuous, and characteristic.
Input coordinates are standardized with training statistics. We use
$\lambda=1$ and AdamW for 8{,}000 steps with batch size 512, and compare
agreement alone and ordinary-mean prediction
(Eq.~\ref{eq:anchor-mean-objective}) under the same encoder and budget.
For each factor, a Gaussian kernel ridge readout is fitted on 3{,}000
training codes, with bandwidth, ridge, and input whitening selected on
validation data, and scored by the test $R^2$ in
Eq.~\ref{eq:recovery-r2}. Content $R^2$ applies the same readout to
predict $c$. The same readout applied to the content-free positions
$p_{1:T}$ gives the information available in each window, reported as the
oracle. Because the readout standardizes codes, we also report the ratio
of the code's standard deviation after training to that at
initialization; values near zero indicate collapse toward a constant.

\begin{table}[!htbp]
 \centering
 \caption{\textbf{Action recovery from appearance-resampled pairs.}
 Held-out $R^2$ of the midpoint position $p$, velocity $v$, and
 acceleration $u$, averaged over three seeds. The oracle reads the
 content-free positions of a $T$-frame window. Content is the $R^2$ of
 predicting $c$ from the code of our objective. For the baselines, Avg.\
 is the mean $R^2$ over $(p,v,u)$ and Scale is the final-to-initial ratio
 of the code's standard deviation. Standard deviations across seeds are
 at most $0.03$ for the oracle and our objective and at most $0.14$ for
 the baselines.}
 \label{tab:numerical-window}
 \small
 \setlength{\tabcolsep}{3.4pt}
 \begin{tabular}{@{}lcccccccccccc@{}}
 \toprule
 & & \multicolumn{3}{c}{Oracle}
 & \multicolumn{4}{c}{Agreement + kernel prediction}
 & \multicolumn{2}{c}{Mean prediction}
 & \multicolumn{2}{c}{Agreement only}\\
 \cmidrule(lr){3-5}\cmidrule(lr){6-9}\cmidrule(lr){10-11}\cmidrule(lr){12-13}
 Mixing & $T$ & $p$ & $v$ & $u$ & $p$ & $v$ & $u$ & Content
 & Avg. & Scale & Avg. & Scale\\
 \midrule
 Shear & 1 & 1.00 & 0.00 & 0.00 & 1.00 & 0.00 & 0.00 & 0.53 & 0.33 & 10 & 0.33 & 0.025\\
       & 2 & 0.81 & 1.00 & 0.25 & 0.81 & 1.00 & 0.25 & 0.01 & 0.69 & 13 & 0.66 & 0.070\\
       & 3 & 1.00 & 1.00 & 1.00 & 1.00 & 1.00 & 0.97 & 0.00 & 1.00 & 12 & 0.87 & 0.067\\
       & 4 & 1.00 & 1.00 & 1.00 & 1.00 & 1.00 & 0.97 & 0.00 & 1.00 & 14 & 0.85 & 0.078\\
       & 6 & 1.00 & 1.00 & 1.00 & 1.00 & 1.00 & 0.98 & 0.00 & 1.00 & 9.4 & 0.87 & 0.038\\
       & 8 & 1.00 & 1.00 & 1.00 & 1.00 & 1.00 & 0.94 & 0.00 & 1.00 & 13 & 0.84 & 0.053\\
 \midrule
 Multiplicative & 1 & 1.00 & 0.00 & 0.00 & 1.00 & 0.00 & 0.00 & 0.19 & 0.33 & 0.002 & 0.33 & 0.001\\
       & 2 & 0.81 & 1.00 & 0.24 & 0.81 & 1.00 & 0.24 & 0.00 & 0.59 & 0.003 & 0.59 & 0.003\\
       & 3 & 1.00 & 1.00 & 1.00 & 1.00 & 1.00 & 1.00 & 0.00 & 0.73 & 0.004 & 0.64 & 0.003\\
       & 4 & 1.00 & 1.00 & 1.00 & 1.00 & 1.00 & 1.00 & 0.00 & 0.68 & 0.004 & 0.63 & 0.004\\
       & 6 & 1.00 & 1.00 & 1.00 & 1.00 & 1.00 & 1.00 & 0.00 & 0.58 & 0.003 & 0.65 & 0.003\\
       & 8 & 1.00 & 1.00 & 1.00 & 1.00 & 1.00 & 1.00 & 0.00 & 0.61 & 0.004 & 0.59 & 0.004\\
 \bottomrule
 \end{tabular}
\end{table}

\paragraph{Observation window.}
Table~\ref{tab:numerical-window} follows
Proposition~\ref{prop:nonlinear-temporal-example}. With $T\geq3$, agreement
plus kernel prediction recovers all three action parameters under both
mixing maps, with content $R^2$ of $0.00$; at a fixed action, less than
$0.1\%$ of the code variance changes with content. With $T=2$, recovery
equals the oracle, whose loss of $u$ reflects the unobservable direction
rather than learning error. With $T=1$ only $p$ is observable; the code has
more coordinates than the window identifies, and its low-variance
coordinates retain content. On the shear map, the conditional feature mean
$(q_t^2,\,q_t+q_t^4+1/3)$ with $q_t=p_t/2.5$ is injective, and ordinary-mean
prediction also identifies the action, consistent with
Lemma~\ref{lem:anchor-learning}. On the multiplicative map, whose
conditional mean vanishes, ordinary-mean prediction and agreement alone
shrink the code to less than $1\%$ of its initial scale. The readout still
extracts part of the action from the remaining variation, but 4--6\% of
that variation changes with content at a fixed action. Agreement alone
also shrinks the code on the shear map and recovers the action less
completely. Prediction without agreement recovered the action without
content for $T\geq3$ in these runs ($R^2\geq0.998$, content $R^2\leq0.00$),
although its population minimizers need not exclude content.

With four-dimensional content passed through a random three-layer
invertible LeakyReLU network, the oracle remains $1.00$ for $T\geq3$, but
recovery by our objective at this training budget is incomplete and
seed-dependent: $R^2$ is $0.82$--$0.87$ for $p$, $0.95$--$0.97$ for $v$, and
$0.36$--$0.47$ for $u$, with seed standard deviations up to $0.49$ and
content $R^2$ of $0.05$--$0.06$. This gap reflects the finite-data approximation
rather than the observation window.

\paragraph{Mean counterexample and premise departures.}
For Eq.~\ref{eq:nonlinear-mean-counterexample} with a one-dimensional code,
our objective recovers $a$ with $R^2=1.000$ and content $R^2=-0.001$.
Ordinary-mean prediction and agreement alone shrink the code to $0.3\%$
and $0.1\%$ of its initial scale, and their readouts reach
$R^2=0.65\pm0.23$ and $0.10\pm0.15$. On the shear map with $T=4$, sharing
content between the two views, which violates
Condition~\ref{cond:anchor-identification}(ii), leaves content in the code:
content $R^2$ is $0.997$, and $81\%$ of the code variance changes with
content at a fixed action. Averaging per-frame features over time instead
of reading the ordered video gives $R^2=0.00$ for $v$, because reversal maps
$(p,v,u)$ to $(p,-v,u)$ within the support, and $0.16$ for $u$.
Independent observation noise with standard deviation $0.05$ makes the
observations non-injective and lowers $u$ to about $0.5$. Recovery is stable
for $\lambda\in\{0.1,10\}$, with $R^2$ for $u$ of $0.98$ and $0.93$, and
for an eight-dimensional code, with $R^2=1.00$ for all three parameters and
content $R^2=0.06$.

\paragraph{Additive constructions.}
For the linear construction of Eq.~\ref{eq:temporal-kinematic-example}
with $\mathbf x_t=\mathbf M[n;p_t]$, a random nonsingular $\mathbf M$,
times $t=1,\ldots,T$, and uniform $\mathbf w$ and $n$, the covariance
projection of Proposition~\ref{prop:structural-anchor-recovery} is
estimated from the 20{,}000 training pairs alone. Its dimension is
$\min(T,3)$, as predicted, and paired codes agree to floating-point
precision, with a within-action variance fraction below $10^{-14}$.
It recovers $\mathbf w$ with $R^2=1.000$ for $T\in\{3,4,6\}$, $0.704$ for
$T=2$, and $0.335$ for $T=1$, with content $R^2\leq0.001$ throughout.
For the squared dynamics of Eq.~\ref{eq:nonlinear-action-construction},
the projection recovers $(p_0,u)$ with $R^2=0.995$ for $T=3$ and $0.525$
for $T=2$. With independent observation noise, the paired-difference
covariance has full rank: at noise standard deviation $0.01$, no exactly
invariant direction remains in the linear construction, as anticipated in
Appendix~\ref{subsec:structural-anchor-proof}. Taking the three leading
generalized eigenvectors of $(\boldsymbol\Sigma_\Phi,\boldsymbol\Sigma_\Delta)$
instead gives $R^2=0.996$ at noise $0.01$ and $0.849$ at noise $0.1$ for
$T=3$. Learned ordinary-mean prediction, analyzed in
Corollary~\ref{cor:additive-anchor-learning}, recovers $\mathbf w$ with
$R^2=1.000$ for $T\geq3$ on the linear construction; our objective reaches
$0.80\pm0.15$ for $T=3$ and $1.00$ for $T=6$. On the squared dynamics,
whose features span three orders of magnitude, both learned targets reach
$R^2\approx0.82$ for $T=3$.

Across these settings, recovery follows the stated conditions. It is
complete once the window is observable and appearance is resampled
independently, it matches the oracle when the window is too short, and it
degrades when paired views share content or temporal order is discarded.
The kernel target retains actions whose conditional feature mean vanishes,
where ordinary-mean prediction collapses.
